\section{Datasets and Experimental Protocol}

\subsection{Dataset roles and composition}

The project uses datasets according to what their evidence can support. D-Link
Camera 5 supplies the reference-study workload and has no attack ground truth in
the repository. The two UNSW datasets contain mixed-device captures and official
attack annotations; preprocessing isolates the selected device by MAC address.
Source files remain unchanged, and label preparation produces separate,
time-aligned JSON records.

Complete captures, not random rows, are assigned to each role. Table
\ref{tab:dataset-composition} reports the number of processed ten-minute windows
currently stored on \texttt{master}.

\begin{table}[H]
\centering
\small
\begin{tabularx}{\textwidth}{>{\raggedright\arraybackslash}p{0.22\textwidth}
>{\raggedright\arraybackslash}p{0.15\textwidth} X rr}
\toprule
\textbf{Dataset} & \textbf{Role} & \textbf{Capture dates} & \textbf{Windows}
& \textbf{Attack windows} \\
\midrule
D-Link Camera 5 & Fit & Oct. 8--13, 2020 & 846 & -- \\
{} & Calibration & Oct. 14--15, 2020 & 288 & -- \\
{} & Development & Oct. 16 and 19, 2020 & 270 & -- \\
{} & Final & Oct. 21--22, 2020 & 288 & -- \\
\midrule
UNSW Chromecast & Fit & Oct. 10--15, 2018 & 867 & 0 \\
{} & Calibration & Oct. 16--19, 2018 & 576 & 0 \\
{} & Development & Oct. 20--23, 2018 & 474 & 32 \\
{} & Final & Oct. 25--27, 2018 & 432 & 6 \\
\midrule
UNSW Samsung camera & Fit & May 28--30, 2018 & 432 & 0 \\
{} & Calibration & May 31, 2018 & 81 & 0 \\
{} & Development & June 1--4, 2018 & 556 & 70 \\
{} & Final & June 5--8, 2018 & 576 & 12 \\
\bottomrule
\end{tabularx}
\caption{Stored window composition by experimental role. Dashes indicate that
attack labels are unavailable.}
\label{tab:dataset-composition}
\end{table}

The Chromecast October 24 source is excluded because the publisher's capture is
truncated; six annotations on that date are consequently outside the retained
dataset.

\subsection{Source preparation and annotation alignment}

The retained D-Link raw archive contains 64 tracked captures, but its research
configuration selects twelve. Only the selected 1,692 processed windows enter
the four declared research roles. A large raw directory
does not imply that every capture was used for model fitting.

The prepared Chromecast and Samsung views use symbolic links to downloaded
sources rather than duplicating PCAP bytes. The download scripts keep
device-specific benign and mixed staging, resume transfers, and reproduce the
declared date selection. The preparation script exposes the flat source view
expected by the configured processor. Source bytes remain separate from
filtered device windows.

\path{prepare_labels.py} selects the official device annotation file and
derives window bounds from packet timestamps, not just the nominal filename.
A label is positive when
$a_{\mathrm{start}}<w_{\mathrm{end}}$ and
$a_{\mathrm{end}}>w_{\mathrm{start}}$.
The union of overlapping annotation intervals determines overlap seconds,
preventing duplicate attack intervals from exceeding the window's duration.
Labels and processed counts remain separate JSON files joined only offline.

The Chromecast record contains 27 official attack intervals, of which 21
match retained captures. It yields 2,349 labelled windows and 38 positives:
32 in development and six in final data. October 21 ends after about seven
hours and contributes only 42 windows. The removed October 24 source accounts
for the six unmatched intervals. Its publisher-side truncation is a recorded
data-selection decision, not a detector failure.

Samsung's 42 intervals all match retained captures. Its 1,645 labelled windows
include 82 positives, with 70 in development and twelve in final data.
May 31 contains 81 windows and June 1 contains 124, rather than a full 144.
The labels therefore measure the observed recording coverage; a nominal day
does not automatically supply a full day of evidence.

Direct inspection finds 187 observed-silent Chromecast development windows,
30 silent Samsung development windows, and 19 silent Samsung final windows.
No missing-state rows are stored in these three configured collections.
That absence does not prove gap-free source coverage; explicit synthetic
discontinuities and the edge fixture exercise the missing state separately.

\subsection{Chronology and leakage policy}

Model fitting and calibration use different earlier captures. The development
stream comes later, and declared final captures come later again. No label is
used to fit alpha, choose a calibration threshold, or decide whether an adaptive
window enters a buffer. Development labels are allowed only to compare completed
runs and diagnose failures.

This capture-level split is preferable to random window sampling, but it is not
the end of the temporal problem. Adjacent scores remain correlated, one boundary
provides only one version of history, and the master validator checks dates in
capture identifiers rather than every processed timestamp. The separate temporal
branch addresses those limitations in Section~\ref{sec:temporal-wip}.
